\documentclass[book.tex]{subfiles}
\usepackage{float}
\usepackage{hyperref}
\usepackage{listings}
\begin{document}
\chapter{Introduction to Scientific Machine Learning}
\label{ch:intro-molecular-ml}
While applying machine learning algorithms to physical systems, there typically are additional constraints in the form of respecting any given laws of physics often described by general nonlinear partial differential equations. In this chapter, we will introduce approaches to incorporate such additional constraints into machine learning algorithms. We will also discuss approaches to discover and learn the underlying partial differential equations for systems where the underlying physics can emerge from collected data. 

This chapter will cover the three broad ways in which one can impose the known physics and symmetries in the physical systems. Often, a combination of the three broad verticals are employed. 

\begin{itemize}
    \item Chapter Part-1: Hard-code the known governing equations and boundary conditions in the \textbf{loss function} while training the physics-driven model, which comprises of approaches that fall under the category of \textbf{physics-informed models} 
    \item Chapter Part-2: Build in the physics-based equations and constraints \textbf{within the model architecture} so that intermediate variables are physically motivated
    \item Chapter Part-3: Impose known symmetries of the system that need to be obeyed \textbf{within the featurizer(s)} before feeding into the model
\end{itemize}
\section{Physics-informed Machine Learning}
Let's begin with discussing physics-informed machine learning using neural networks. The approach by Raissi et al.\cite{raissi2019physics} employs deep neural networks to leverage their well known capability as universal function approximators. This facilitatest the ability to directly tackle nonlinear problems without the need for committing to any prior assumptions, linearization, or local time-stepping. In addition we can exploit recent developments in automatic differentiation – one of the most useful but perhaps under-utilized techniques in scientific computing – to differentiate neural networks with respect to their input coordinates and model parameters to obtain physics-informed neural networks. Such neural networks are constrained to respect any symmetries, invariances, or conservation principles originating from the physical laws that govern the observed data, as modeled by general time-dependent and nonlinear partial differential equations. This simple yet powerful construction allows us to tackle a wide range of problems in computational science and introduces a potentially transformative technology leading to the development of new data-efficient and physics-informed learning machines, new classes of numerical solvers for partial differential equations, as well as new data-driven approaches for model inversion and systems identification. The approach sets the foundations for a new paradigm in modeling and computation that enriches deep learning with the longstanding developments in mathematical physics.

\subsection{Problem Formulation}
A parametrized and nonlinear partial differential equations of the following general form are considered 
\begin{equation}
    u_t + \mathcal{N}[u;\lambda]=0, x \in \Omega, t\in [0,T],
\end{equation}
where $u(t, x)$ denotes the hidden solution, $\mathcal{N[\cdot;\lambda]}$ is a nonlinear operator parametrized by $\lambda$, and $\Omega$ is a subset of $\mathbb{R}^D$. This formulation applies to a wide range of problems in physical systems including conservation laws, diffusion processes, kinetic equations and physical systems.

\subsection{Data-driven solutions of partial differential equations}
Let's begin with solutions to formulations of the general form 
\begin{equation}
    u_t + \mathcal{N}[u]=0, x \in \Omega, t\in [0,T],
    \label{eq_generalform}
\end{equation}
where $u(t, x)$ denotes the hidden solution, $\mathcal{N[\cdot]}$ is a nonlinear operator, and $\Omega$ is a subset of $\mathbb{R}^D$.

The approach discussed below pertains to an algorithm that can be used on \textbf{continuous time models}.

Let's define $f(t, x)$ to be given by the equation on the left hand side,

\begin{equation}
    f:=u_t + \mathcal{N}[u]
    \label{eq_contTime}
\end{equation}

and the approach is to approximate $u(t,x)$ by a deep neural network. This assumption along with equation \ref{eq_contTime} gives rise to the physics-informed neural network $f(t,x)$. The network can be derived by applying the chain rule for differentiating compositions of functions using automatic differentiation, and has the same parameters as the network representing $u(t, x)$, albeit with different activation functions due to the action of the differential operator $\mathcal{N}$. The shared parameters between the neural networks $u(t, x)$ and $f(t, x)$ can be learned by minimizing the mean-squared error loss

\begin{equation}
    MSE=MSE_u + MSE_f 
\end{equation}



where 
\begin{equation}
    MSE_u=\frac{1}{N_u} \sum_{i=1}^{N_u}|u(t_u^i,x_u^i)-u^i|^2
\end{equation}

and 
\begin{equation}
    MSE_f=\frac{1}{N_f} \sum_{i=1}^{N_f}|f(t_f^i,x_f^i)|^2
\end{equation}

In the problem formulation, the initial and boundary training data on $u(t, x)$ are denoted by $\{{t^i_u, x^i_u, u^i}\}^{N_u}_{i=1}$, and $\{{t^i_f, x^i_f}\}^{N_f}_{i=1}$ specify the collocation points for $f(t, x)$. 
The loss function $MSE_u$ corresponds to the initial and boundary data while $MSE_f$ enforces the structure imposed by the general equation (equation \ref{eq_generalform}) at a finite set of collocation points. 

\begin{example}[Schr{\"o}dinger equation]

This example aims to highlight the ability of The method to handle periodic boundary conditions, complex-valued solutions, as well as different types of nonlinearities in the governing partial differential equations. The nonlinear Schr$\mathrm{\ddot{o}}$dinger equation along with periodic boundary conditions is given by


\begin{align}
\begin{split}
i h_t + 0.5 h_{xx} + |h|^2 h = 0,\ \ \ x \in [-5, 5],\ \ \ t \in [0, \pi/2],\\
h(0,x) = 2\ \mathrm{sech}(x),\\
h(t,-5) = h(t, 5),\\
h_x(t,-5) = h_x(t, 5),
\end{split}
\end{align}


where $h(t,x)$ is the complex-valued solution. Let us define $f(t,x)$ to be given by

\begin{equation}
    f := i h_t + 0.5 h_{xx} + |h|^2 h,    
\end{equation}

and proceed by placing a complex-valued neural network prior on $h(t,x)$. In fact, if $u$ denotes the real part of $h$ and $v$ is the imaginary part, we are placing a multi-out neural network prior on $h(t,x) = \begin{bmatrix}
u(t,x) & v(t,x)
\end{bmatrix}$. This will result in the complex-valued (multi-output) physics informed neural network $f(t,x)$. The shared parameters of the neural networks $h(t,x)$ and $f(t,x)$ can be learned by minimizing the mean squared error loss

\begin{equation}
    MSE = MSE_0 + MSE_b + MSE_f    
\end{equation}


where

\begin{equation}
MSE_0 = \frac{1}{N_0}\sum_{i=1}^{N_0} |h(0,x_0^i) - h^i_0|^2    
\end{equation}


\begin{equation}
MSE_b = \frac{1}{N_b}\sum_{i=1}^{N_b} \left(|h^i(t^i_b,-5) - h^i(t^i_b,5)|^2 + |h^i_x(t^i_b,-5) - h^i_x(t^i_b,5)|^2\right)    
\end{equation}


and

\begin{equation}
MSE_f = \frac{1}{N_f}\sum_{i=1}^{N_f}|f(t_f^i,x_f^i)|^2.    
\end{equation}

Here, $\{x_0^i, h^i_0\}_{i=1}^{N_0}$ denotes the initial data, $\{t^i_b\}_{i=1}^{N_b}$ corresponds to the collocation points on the boundary, and $\{t_f^i,x_f^i\}_{i=1}^{N_f}$ represents the collocation points on $f(t,x)$. Consequently, $MSE_0$ corresponds to the loss on the initial data, $MSE_b$ enforces the periodic boundary conditions, and $MSE_f$ penalizes the Schr$\mathrm{\ddot{o}}$dinger equation not being satisfied on the collocation points.

\begin{figure}
    \centering
    \includegraphics{figures/NLS_bookl.png}
    \caption{Schrödinger equation: \textbf{Top:} Predicted
solution along with the initial and boundary training data. In addition
we are using 20,000 collocation points generated using a Latin Hypercube
Sampling strategy. \textbf{Bottom:} Comparison of the predicted and
exact solutions corresponding to the three temporal snapshots depicted
by the dashed vertical lines in the top panel.}
\end{figure}




\end{example}
% The following figure summarizes the results of The experiment.


One potential limitation of the continuous time neural network models considered so far stems from the need to use a large number of collocation points $N_f$ in order to enforce physics informed constraints in the entire spatio-temporal domain. Although this poses no significant issues for problems in one or two spatial dimensions, it may introduce a severe bottleneck in higher dimensional problems, as the total number of collocation points needed to globally enforce a physics informed constrain (i.e., in this case a partial differential equation) will increase exponentially. In the next section, we put forth a different approach that circumvents the need for collocation points by introducing a more structured neural network representation leveraging the classical \href{https://en.wikipedia.org/wiki/Runge–Kutta_methods}{Runge-Kutta} time-stepping schemes.


%%%%%%%%%


The approach discussed below pertains to a methodology that can be used on \textbf{discrete-time models}. Let's first revist the classical Runge-Kutta method (4 steps) before discussing the general case of $q$ steps.  

\paragraph{The Classical Runge-Kutta
Method}\label{the-classical-runge-kutta-method}

Let an initial value problem be specified as follows:

$\frac{dy}{dt}=f(t,y)$,$\quad y(t_{0})=y_{0}$

Now pick a step-size h \textgreater{} 0 and define

$\begin{aligned}y_{n+1}&=y_{n}+{\frac {1}{6}}h\left(k_{1}+2k_{2}+2k_{3}+k_{4}\right),\\t_{n+1}&=t_{n}+h\\\end{aligned}$

$
\begin{aligned}k_{1}&=\ f(t_{n},y_{n}),\\k_{2}&=\ f\left(t_{n}+{\frac {h}{2}},y_{n}+h{\frac {k_{1}}{2}}\right),\\k_{3}&=\ f\left(t_{n}+{\frac {h}{2}},y_{n}+h{\frac {k_{2}}{2}}\right),\\k_{4}&=\ f\left(t_{n}+h,y_{n}+hk_{3}\right).\end{aligned}
$
\paragraph{The General Form of the Runge-Kutta
Method}\label{the-general-runge-kutta-method}
The general form of Runge-Kutta methods with $q$ stages is applied to the governing equation to arrive at the following:

\begin{align}
    \begin{split}
        u^{n+c_i}=u^n-\Delta t\sum_{i=1}^{q}a_{ij}\mathcal{N}[u^{n+c_j}], i=1,\cdot\cdot\cdot,q,\\
        u^{n+1}=u^n-\Delta t\sum_{i=1}^{q}b_{j}\mathcal{N}[u^{n+c_j}].
    \end{split}
\end{align}

In the formulation, $u^{n+c_j}(x)=u(t^n+c_j\Delta t,x)$ for $j=1,\cdot\cdot\cdot,q$. Depending on the choice of the parameters $\{a_{ij},b_j,c_j\}$.

The above equations can also be expressed as

\begin{align}
    \begin{split}
        u^n=u^n_i, i=1,\cdot\cdot\cdot,q,\\
        u^n=u^n_{q+1}
    \end{split}
\end{align}

where
\begin{align}
    \begin{split}
        u_i^n:=u^{n+c_i}+\Delta t\sum_{j=1}^{q}a_{ij}\mathcal{N}[u^{n+c_j}], i=1,\cdot\cdot\cdot,q,\\
          u^{n}_{q+1}=u^{n+1}+\Delta t\sum_{i=1}^{q}b_{j}\mathcal{N}[u^{n+c_j}].
    \end{split}
\end{align}

We can place a multi-output neural network prior on
\begin{equation}
\label{eq_discrete}
    [u^{n+c_1}(x),\cdot\cdot\cdot,u^{n+c_q}(x),u^{n+1}(x)]
\end{equation}
The assumption along with equations \ref{eq_discrete} result in a physics-informed neural network that takes $x$ as an input and gives as the output 
\begin{equation}
    [u^{n}_1(x),\cdot\cdot\cdot,u^{n}_q(x),u_{q+1}^{n+1}(x)]
\end{equation}

\begin{example}{Burgers' Equation}\\
The Burgers' equation in one space dimension along with Dirichlet boundary conditions reads as follow:

\[
\begin{array}{l}
u_t + u u_x - (0.01/\pi) u_{xx} = 0,\ \ \ x \in [-1,1],\ \ \ t \in [0,1],\\
u(0,x) = -\sin(\pi x),\\
u(t,-1) = u(t,1) = 0.
\end{array}
\]

Let us define $f(t,x)$ to be given by

\[
f := u_t + u u_x - (0.01/\pi) u_{xx},
\]

and proceed by approximating $u(t,x)$ by a deep neural network.

Correspondingly, the physics informed neural network $f(t,x)$ takes the
form

Loss function: $MSE = MSE_u + MSE_f,$

where

$ MSE\_u = \frac{1}{N_u}\sum\emph{\{i=1\}\^{}\{N}u\}
\textbar{}u(t\textsuperscript{i\_u,x\_u}i) -
u\textsuperscript{i\textbar{}}2$ and
$MSE_f = \frac{1}{N_f}\sum_{i=1}^{N_f}|f(t_f^i,x_f^i)|^2$


The following figure summarizes the results for the data-driven solution
of the Burgers' equation.

\begin{figure}
    \centering
    \includegraphics{figures/Burgers_CT_inference_book.png}
    \caption{\emph{Top:} Predicted solution along with the initial and boundary training data. In addition we are using 10,000 collocation points generated using a Latin Hypercube Sampling strategy. \emph{Bottom:} Comparison of the predicted and exact solutions corresponding to the three temporal snapshots depicted by the white vertical lines in the top panel. Model training took approximately 60 seconds on a single NVIDIA Titan X GPU card.}
    \label{fig:Burgers's}
\end{figure}
\end{example}

\begin{example}[Allen-Cahn Equations]
 Below we show an example for the Allen-Cahn equation that appears in systems with phase separations, e.g., battery electrode-electrolyte interfaces. The example is used to demonstrate the ability of the proposed discrete time
 models to handle different types of nonlinearity in the governing partial differential equation. Let's consider the Allen–Cahn equation along with periodic boundary conditions
 
 \begin{align}
     u_t-0.0001u_{xx}+5u^3-5u^3-5u &= 0, x\in[-1,1], t \in [0,1]\\
     u(0,x) &=x^2 cos(\pi x)\\
     u(t,-1) &=u(t,1)\\
     u_x(t,-1) &=u_x(t,1)\\
 \end{align}
 
 For the Allen-Cahn equation, the non-linear operator is given by
 
 \begin{equation}
     \mathcal{N}[u^{n+c_j}]=-0.0001u^{n+c_j}_{xx}+5(u^{n+c_j})^3-5u^{n+c_j}
 \end{equation}
 
 and the shared parameters of the neural networks can be learnt by minimizing the loss function as discussed earlier,
 
 \begin{equation}
     SSE=SSE_n+SSE_b
 \end{equation}
 
 where
 \begin{equation}
     SSE_n=\sum_{j=1}^{q+1}\sum_{i=1}^{N_n}|u^n_j(x^{n,i})=u^{n,i}|^2
 \end{equation}
 
 and
 \begin{align}
 \begin{split}
     SSE_b=\sum_{i=1}^q|u^{n+c_i}(-1)-u^{n+c_i}(1)|^2+|u^{n+1}(-1)-u^{n+1}(1)|^2\\
     +\sum_{i=1}^q|u_x^{n+c_i}(-1)-u_x^{n+c_i}(1)|^2+|u_x^{n+1}(-1)-u_x^{n+1}(1)|^2
 \end{split}
 \end{align}

\begin{figure}[h]
     \centering
    \includegraphics{figures/allenCahn_PINN.png}
    \caption{Allen-Cahn Equation: Discrete Time Domain}
    \label{fig:allencahn}
\end{figure}

\end{example}

\section{Data-Driven Discovery of Nonlinear Partial Differential Equations}

We shift our attention to the problem of data-driven discovery of partial differential equations. To this end, let us consider parametrized and nonlinear partial differential equations of the general form
\begin{equation}
u_t + \mathcal{N}[u;\lambda] = 0,\ x \in \Omega, \ t\in[0,T],    
\end{equation}


where $u(t,x)$ denotes the latent (hidden) solution, $\mathcal{N}[\cdot;\lambda]$ is a nonlinear operator parametrized by $\lambda$, and $\Omega$ is a subset of $\mathbb{R}^D$. Now, the problem of data-driven discovery of partial differential equations poses the following question: given a small set of scattered and potentially noisy observations of the hidden state $u(t,x)$ of a system, what are the parameters $\lambda$ that best describe the observed data?

In what follows, we will provide an overview of the two main approaches to tackle this problem, namely continuous time and discrete time models, as well as a series of results and systematic studies for a diverse collection of benchmarks. In the first approach, we will assume availability of scattered and potential noisy measurements across the entire spatio-temporal domain. In the latter, we will try to infer the unknown parameters $\lambda$ from only two data snapshots taken at distinct time instants

\subsection{Continuous Time Models}
We define $f(t,x)$ to be given by

\begin{equation}
f := u_t + \mathcal{N}[u;\lambda],\label{eq:PDE_RHS}    
\end{equation}



and proceed by approximating $u(t,x)$ by a deep neural network. This assumption results in a \href{https://arxiv.org/abs/1711.10566}{physics informed neural network} $f(t,x)$. This network can be derived by the calculus on computational graphs: \href{http://colah.github.io/posts/2015-08-Backprop/}{Backpropagation}. It is worth highlighting that the parameters of the differential operator $\lambda$ turn into parameters of the physics informed neural network $f(t,x)$.

\begin{example}{Navier-Stokes Equation}

The next example involves a realistic scenario of incompressible fluid flow as described by the Navier-Stokes equations. Navier-Stokes equations describe the physics of many phenomena of scientific and engineering interest. They may be used to model the weather, ocean currents, water flow in a pipe and air flow around a wing. The Navier-Stokes equations in their full and simplified forms help with the design of aircraft and cars, the study of blood flow, the design of power stations, the analysis of the dispersion of pollutants, and many other applications. Let us consider the Navier-Stokes equations in two dimensions (2D) given explicitly by


\begin{align}
\begin{split}
u_t + \lambda_1 (u u_x + v u_y) = -p_x + \lambda_2(u_{xx} + u_{yy}),\\
v_t + \lambda_1 (u v_x + v v_y) = -p_y + \lambda_2(v_{xx} + v_{yy}),    
\end{split}
\end{align}


where $u(t, x, y)$ denotes the $x$-component of the velocity field, $v(t, x, y)$ the $y$-component, and $p(t, x, y)$ the pressure. Here, $\lambda = (\lambda_1, \lambda_2)$ are the unknown parameters. Solutions to the Navier-Stokes equations are searched in the set of divergence-free functions; i.e.,

\begin{equation}
u_x + v_y = 0.    
\end{equation}


This extra equation is the continuity equation for incompressible fluids that describes the conservation of mass of the fluid. We make the assumption that

\begin{equation}
u = \psi_y,\ \ \ v = -\psi_x,    
\end{equation}


for some latent function $\psi(t,x,y)$. Under this assumption, the continuity equation will be automatically satisfied. Given noisy measurements

\begin{equation}
\{t^i, x^i, y^i, u^i, v^i\}_{i=1}^{N}    
\end{equation}


of the velocity field, we are interested in learning the parameters $\lambda$ as well as the pressure $p(t,x,y)$. We define $f(t,x,y)$ and $g(t,x,y)$ to be given by

\begin{align}
f &:= u_t + \lambda_1 (u u_x + v u_y) + p_x - \lambda_2(u_{xx} + u_{yy})\\
g &:= v_t + \lambda_1 (u v_x + v v_y) + p_y - \lambda_2(v_{xx} + v_{yy})        
\end{align}

\noindent and proceed by jointly approximating \[\begin{bmatrix}
\psi(t,x,y) & p(t,x,y)
\end{bmatrix} \] 
using a single neural network with two outputs. This prior assumption results into a physics informed neural network \[\begin{bmatrix}
f(t,x,y) & g(t,x,y)
\end{bmatrix} \]
The parameters $\lambda$ of the Navier-Stokes operator as well as the parameters of the neural networks $\begin{bmatrix}
\psi(t,x,y) & p(t,x,y)
\end{bmatrix}$ and $\begin{bmatrix}
f(t,x,y) & g(t,x,y)
\end{bmatrix}$ can be trained by minimizing the mean squared error loss


\begin{align}
    \begin{split}
    MSE :=& \frac{1}{N}\sum_{i=1}^{N} \left(|u(t^i,x^i,y^i) - u^i|^2 + |v(t^i,x^i,y^i) - v^i|^2\right) \\
    +& \frac{1}{N}\sum_{i=1}^{N} \left(|f(t^i,x^i,y^i)|^2 + |g(t^i,x^i,y^i)|^2\right)
    \end{split}
\end{align}


\begin{figure}
    \centering
    \includegraphics{figures/NavierStokes_data_book.png}
    \caption{Navier-Stokes equation: \textbf{Top:} Incompressible flow and dynamic vortex shedding past a circular cylinder at Re$=100$. The spatio-temporal training data correspond to the depicted rectangular region in the cylinder wake. \textbf{Bottom:} Locations of
training data-points for the the stream-wise and transverse velocity
components.
}
\end{figure}

\begin{figure}
    \centering
    \includegraphics{figures/NavierStokes_prediction_book.png}
    \caption{Navier-Stokes equation: \textbf{Top:} Predicted
versus exact instantaneous pressure field at a representative time
instant. By definition, the pressure can be recovered up to a constant,
hence justifying the different magnitude between the two plots. This
remarkable qualitative agreement highlights the ability of
physics-informed neural networks to identify the entire pressure field,
despite the fact that no data on the pressure are used during model
training. \textbf{Bottom:} Correct partial differential equation along
with the identified one.}
\end{figure}

\end{example}

The approach so far assumes availability of scattered data throughout the entire spatio-temporal domain. However, in many cases of practical interest, one may only be able to observe the system at distinct time instants. In the next section, we introduce a different approach that tackles the data-driven discovery problem using only two data snapshots. We will see how, by leveraging the classical Runge-Kutta time-stepping schemes, one can construct discrete time physics informed neural networks that can retain high predictive accuracy even when the temporal gap between the data snapshots is very large.


\subsection{Discrete Time Models}

We begin by employing the general form of Runge-Kutta methods with $q$ stages and obtain

\begin{align}
u^{n+c_i} &= u^n - \Delta t \sum_{j=1}^q a_{ij} \mathcal{N}[u^{n+c_j};\lambda], \ \ i=1,\ldots,q\\
u^{n+1} &= u^{n} - \Delta t \sum_{j=1}^q b_j \mathcal{N}[u^{n+c_j};\lambda].        
\end{align}


Here, $u^{n+c_j}(x) = u(t^n + c_j \Delta t, x)$ for $j=1, \ldots, q$. This general form encapsulates both implicit and explicit time-stepping schemes, depending on the choice of the parameters $\{a_{ij},b_j,c_j\}$. The above equations can be equivalently expressed as

\begin{align}
u^{n} &= u^n_i, \ \ i=1,\ldots,q\\
u^{n+1} &= u^{n+1}_{i}, \ \ i=1,\ldots,q
\end{align}

where

\begin{align}
u^n_i &:= u^{n+c_i} + \Delta t \sum_{j=1}^q a_{ij} \mathcal{N}[u^{n+c_j};\lambda], \ \ i=1,\ldots,q,\\
u^{n+1}_{i} &:= u^{n+c_i} + \Delta t \sum_{j=1}^q (a_{ij} - b_j) \mathcal{N}[u^{n+c_j};\lambda], \ \ i=1,\ldots,q.
\end{align}



We proceed by placing a multi-output neural network prior on

\begin{equation}
u^{n+c_1}(x), \ldots, u^{n+c_q}(x)    
\end{equation}


This prior assumption result in two physics informed neural networks

\begin{equation}
u^{n}_1(x), \ldots, u^{n}_q(x), u^{n}_{q+1}(x)    
\end{equation}


and

\begin{equation}
u^{n+1}_1(x), \ldots, u^{n+1}_q(x), u^{n+1}_{q+1}(x)    
\end{equation}



Given noisy measurements at two distinct temporal snapshots $\{\mathbf{x}^{n}, \mathbf{u}^{n}\}$ and $\{\mathbf{x}^{n+1}, \mathbf{u}^{n+1}\}$ of the system at times $t^{n}$ and $t^{n+1}$, respectively, the shared parameters of the neural networks along with the parameters $\lambda$ of the differential operator can be trained by minimizing the sum of squared errors

\begin{equation}
SSE = SSE_n + SSE_{n+1},    
\end{equation}



where

\begin{equation}
SSE_n := \sum_{j=1}^q \sum_{i=1}^{N_n} |u^n_j(x^{n,i}) - u^{n,i}|^2,    
\end{equation}



and

\begin{equation}
SSE_{n+1} := \sum_{j=1}^q \sum_{i=1}^{N_{n+1}} |u^{n+1}_j(x^{n+1,i}) - u^{n+1,i}|^2.    
\end{equation}



Here, $\mathbf{x}^n = \left\{x^{n,i}\right\}_{i=1}^{N_n}$, $\mathbf{u}^n = \left\{u^{n,i}\right\}_{i=1}^{N_n}$, $\mathbf{x}^{n+1} = \left\{x^{n+1,i}\right\}_{i=1}^{N_{n+1}}$, and $\mathbf{u}^{n+1} = \left\{u^{n+1,i}\right\}_{i=1}^{N_{n+1}}$.


\begin{example}{Korteweg–de Vries Equation}

The final example aims to highlight the ability of the proposed framework to handle governing partial differential equations involving higher order derivatives. Here, we consider a mathematical model of waves on shallow water surfaces, the Korteweg-de Vries (KdV) equation. The KdV equation reads as

\begin{equation}
u_t + \lambda_1 u u_x + \lambda_2 u_{xxx} = 0,    
\end{equation}

with $(\lambda_1, \lambda_2)$ being the unknown parameters. For the KdV equation, the nonlinear operator is given by

\begin{equation}
\mathcal{N}[u^{n+c_j}] = \lambda_1 u^{n+c_j} u^{n+c_j}_x - \lambda_2 u^{n+c_j}_{xxx}    
\end{equation}


and the shared parameters of the neural networks along with the parameters $\lambda = (\lambda_1, \lambda_2)$ of the KdV equation can be learned by minimizing the sum of squared errors given above.

\begin{figure}
    \centering
    \includegraphics{figures/KdV_book.png}
    \caption{KdV equation: \textbf{Top:} Solution along with
the temporal locations of the two training snapshots. Middle: Training
data and exact solution corresponding to the two temporal snapshots
depicted by the dashed vertical lines in the top panel. \textbf{Bottom:}
Correct partial differential equation along with the identified one.}
    \label{fig:my_label}
\end{figure}

\end{example}

As we have shown in this chapter Physics-informed neural networks are a versatile class of universal function approximators that are capable of encoding any underlying physical laws that govern a given dataset. Physics-informed neural networks enable data-driven algorithms for inferring solutions to general nonlinear partial differential equations. In addition, it enables the construction of computationally efficient physics-informed surrogate models. 

\section{PINNs Implementation and Physika Code}
For the example case of obtaining the solution of the Burgers equation, which is a prototype example of a hyperbolic conservation law, the physika code is discussed below. The simplicity of the implementing the idea can be seen 
To recapitulate, the the Burger’s equation along with Dirichlet boundary conditions reads as follows:

\[
\begin{array}{l}
u_t + u u_x - (0.01/\pi) u_{xx} = 0,\ \ \ x \in [-1,1],\ \ \ t \in [0,1],\\
u(0,x) = -\sin(\pi x),\\
u(t,-1) = u(t,1) = 0.
\end{array}
\]

We define $f(t,x)$ to be given by
\[
f := u_t + u u_x - (0.01/\pi) u_{xx},
\]

$u(t,x)$ is therefore approximated by a deep neural network. 
\begin{lstlisting}
def u(t,x):
    u=neural_net(concat([t,x],1),weights,biases)
    return u
\end{lstlisting}

The physics-informed neural network $f(t, x)$ would be:
\begin{lstlisting}
def f(t,x):
    u=u(t,x)
    u_t=gradients(u,t)[0]
    u_x=gradients(u,x)[0]
    u_xx=gradients(u_x,x)[0]
    f=u_t+(u)(u_x)−(0.01/pi)u_xx
    return f
\end{lstlisting}
The loss function is defined in the following way. A precise description of the loss function is discussed earlier in the chapter for the same application example.
\begin{lstlisting}
def loss:
    MSE_u=(1/N_u) [(u0 - u0_pred) + (v0 - v0_pred) + (u_lb_pred - u_ub_pred)^2 +
                    (v_lb_pred - v_ub_pred)^2 + 
                    (u_x_lb_pred - u_x_ub_pred)^2 + 
                    (v_x_lb_pred - v_x_ub_pred)^2 + 
                    (f_u_pred)^2 + 
                    (f_v_pred)^2)]
\end{lstlisting}
The solution process at a high level can be described as follows:

\begin{lstlisting}
model = PhysicsInformedNN(collocation points, initial conditions, boundary conditions, layers, lower and upper bounds)
model.train(10)
u_pred, v_pred, f_u_pred, f_v_pred = model.predict(X_star)
\end{lstlisting}



















































% \section{PINNs}


% \subsection{Example: (Schrödinger
% Equation)}\label{example-shruxf6dinger-equation}

% \begin{itemize}
% \itemsep1pt\parskip0pt\parsep0pt
% \item
%   handle periodic boundary conditions
% \item
%   complex-valued solutions
% \item
%   different types of nonlinearities in the governing partial
%   differential equations.
% \end{itemize}

% Nonlinear Schrödinger equation:

% \[
% \begin{array}{l}
% i h_t + 0.5 h_{xx} + |h|^2 h = 0,\ \ \ x \in [-5, 5],\ \ \ t \in [0, \pi/2],\\
% h(0,x) = 2\ \text{sech}(x),\\
% h(t,-5) = h(t, 5),\\
% h_x(t,-5) = h_x(t, 5),
% \end{array}
% \]

% where $h(t,x)$ is the complex-valued solution. Let us define $f(t,x)$ to
% be given by

% \[
% f := i h_t + 0.5 h_{xx} + |h|^2 h,
% \]

% Complex-valued neural network prior on $h(t,x)$.

% $u$ denotes the real part of $h$ and $v$ is the imaginary part multi-out
% neural network prior on

% \[h(t,x) = \begin{bmatrix}
% u(t,x) & v(t,x)
% \end{bmatrix}\].

% \subsection{Example: (Schrödinger Equation)
% (contd.)}\label{example-shruxf6dinger-equation-contd.}

% Results in the complex-valued (multi-output) physic informed neural
% network $f(t,x)$. The shared parameters of the neural networks $h(t,x)$
% and $f(t,x)$ can be learned by minimizing the mean squared error loss

% $MSE = MSE_0 + MSE_b + MSE_f,$

% where,

% $MSE_0 = \frac{1}{N_0}\sum_{i=1}^{N_0} |h(0,x_0^i) - h^i_0|^2,$

% $MSE_b = \frac{1}{N_b}\sum_{i=1}^{N_b} \left(|h^i(t^i_b,-5) - h^i(t^i_b,5)|^2 + |h^i_x(t^i_b,-5) - h^i_x(t^i_b,5)|^2\right),$

% and

% $MSE_f = \frac{1}{N_f}\sum_{i=1}^{N_f}|f(t_f^i,x_f^i)|^2.$

% Initial Data: $\{x_0^i, h^i_0\}_{i=1}^{N_0}$

% Collocation points on the boundary: $\{t^i_b\}_{i=1}^{N_b}$

% Collocation points on $f(t,x)$:$\{t_f^i,x_f^i\}_{i=1}^{N_f}$

% $MSE_0$: loss on the initial data, $MSE_b$: enforces the periodic
% boundary conditions, and $MSE_f$ penalizes the Schrödinger equation.

% \subsection{Example: (Schrödinger
% Equation)}\label{example-shruxf6dinger-equation-1}

% The following figure summarizes the results of the experiment.

% \includegraphics{figures/NLS_bookl.png}
% \textgreater{} \emph{Schrödinger equation:} \textbf{Top:} Predicted
% solution along with the initial and boundary training data. In addition
% we are using 20,000 collocation points generated using a Latin Hypercube
% Sampling strategy. \textbf{Bottom:} Comparison of the predicted and
% exact solutions corresponding to the three temporal snapshots depicted
% by the dashed vertical lines in the top panel.

% One potential limitation of the continuous time neural network models
% considered so far, stems from the need to use a large number of
% collocation points $N_f$ in order to enforce physics informed
% constraints in the entire spatio-temporal domain. Although this poses no
% significant issues for problems in one or two spatial dimensions, it may
% introduce a severe bottleneck in higher dimensional problems, as the
% total number of collocation points needed to globally enforce a physics
% informed constrain (i.e., in our case a partial differential equation)
% will increase exponentially. In the next section, we put forth a
% different approach that circumvents the need for collocation points by
% introducing a more structured neural network representation leveraging
% the classical
% \href{https://en.wikipedia.org/wiki/Runge--Kutta_methods}{Runge-Kutta}
% time-stepping schemes.

% \paragraph{Discrete Time Models}\label{discrete-time-models}

% General form of Runge-Kutta methods with $q$ stages and obtain

% \[
% \begin{array}{ll}
% u^{n+c_i} = u^n - \Delta t \sum_{j=1}^q a_{ij} \mathcal{N}[u^{n+c_j}], \ \ i=1,\ldots,q,\\
% u^{n+1} = u^{n} - \Delta t \sum_{j=1}^q b_j \mathcal{N}[u^{n+c_j}].
% \end{array}
% \]

% Here, $u^{n+c_j}(x) = u(t^n + c_j \Delta t, x)$ for $j=1, \ldots, q$.
% This general form encapsulates both implicit and explicit time-stepping
% schemes, depending on the choice of the parameters $\{a_{ij},b_j,c_j\}$.
% The above equations can be equivalently expressed as

% \[
% \begin{array}{ll}
% u^{n} = u^n_i, \ \ i=1,\ldots,q,\\
% u^n = u^n_{q+1},
% \end{array}
% \]

% where

% \[
% \begin{array}{ll}
% u^n_i := u^{n+c_i} + \Delta t \sum_{j=1}^q a_{ij} \mathcal{N}[u^{n+c_j}], \ \ i=1,\ldots,q,\\
% u^n_{q+1} := u^{n+1} + \Delta t \sum_{j=1}^q b_j \mathcal{N}[u^{n+c_j}].
% \end{array}
% \]

% \paragraph{The Classical Runge-Kutta
% Method}\label{the-classical-runge-kutta-method}

% Let an initial value problem be specified as follows:

% $\frac{dy}{dt}=f(t,y)$,$\quad y(t_{0})=y_{0}$

% Now pick a step-size h \textgreater{} 0 and define

% $\begin{aligned}y_{n+1}&=y_{n}+{\frac {1}{6}}h\left(k_{1}+2k_{2}+2k_{3}+k_{4}\right),\\t_{n+1}&=t_{n}+h\\\end{aligned}$

% $
% \begin{aligned}k_{1}&=\ f(t_{n},y_{n}),\\k_{2}&=\ f\left(t_{n}+{\frac {h}{2}},y_{n}+h{\frac {k_{1}}{2}}\right),\\k_{3}&=\ f\left(t_{n}+{\frac {h}{2}},y_{n}+h{\frac {k_{2}}{2}}\right),\\k_{4}&=\ f\left(t_{n}+h,y_{n}+hk_{3}\right).\end{aligned}
% $

% \paragraph{The Classical Runge-Kutta Method
% (Schematic)}\label{the-classical-runge-kutta-method-schematic}

% \begin{figure}[htbp]
% \centering
% \includegraphics{figures/RK_book.png}
% \caption{Slopes at different step locations used by the classical Runge-Kutta method. Source: Wiki}
% \end{figure}

% \paragraph{Discrete Time Models
% (contd.)}\label{discrete-time-models-contd.}

% We proceed by placing a multi-output neural network prior on

% \[
% \begin{bmatrix}
% u^{n+c_1}(x), \ldots, u^{n+c_q}(x), u^{n+1}(x)
% \end{bmatrix}.
% \]

% This prior assumption along with the above equations result in a physics
% informed neural network that takes $x$ as an input and outputs

% \[
% \begin{bmatrix}
% u^n_1(x), \ldots, u^n_q(x), u^n_{q+1}(x)
% \end{bmatrix}.
% \]

% \paragraph{Example: Allen-Cahn
% Equation}\label{example-allen-cahn-equation}

% Allen-Cahn equation along with periodic boundary conditions

% \[
% \begin{array}{l}
% u_t - 0.0001 u_{xx} + 5 u^3 - 5 u = 0, \ \ \ x \in [-1,1], \ \ \ t \in [0,1],\\
% u(0, x) = x^2 \cos(\pi x),\\
% u(t,-1) = u(t,1),\\
% u_x(t,-1) = u_x(t,1).
% \end{array}
% \]

% The Allen-Cahn equation is a well-known equation from the area of
% reaction-diffusion systems. It describes the process of phase separation
% in multi-component alloy systems, including order-disorder transitions.
% For the Allen-Cahn equation, the nonlinear operator is given by

% \[
% \mathcal{N}[u^{n+c_j}] = -0.0001 u^{n+c_j}_{xx} + 5 \left(u^{n+c_j}\right)^3 - 5 u^{n+c_j},
% \]

% and the shared parameters of the neural networks can be learned by
% minimizing the sum of squared errors

% $SSE = SSE_n + SSE_b,$

% \paragraph{Example: Allen-Cahn Equation
% (contd.)}\label{example-allen-cahn-equation-contd.}

% \[SSE_n = \sum_{j=1}^{q+1} \sum_{i=1}^{N_n} |u^n_j(x^{n,i}) - u^{n,i}|^2,\]
% and

% \[
% \begin{array}{rl}
% SSE_b =& \sum_{i=1}^q |u^{n+c_i}(-1) - u^{n+c_i}(1)|^2 + |u^{n+1}(-1) - u^{n+1}(1)|^2 \\
%       +& \sum_{i=1}^q |u_x^{n+c_i}(-1) - u_x^{n+c_i}(1)|^2 + |u_x^{n+1}(-1) - u_x^{n+1}(1)|^2.
% \end{array}
% \]

% Here, $\{x^{n,i}, u^{n,i}\}_{i=1}^{N_n}$ corresponds to the data at time
% $t^n$.

% \paragraph{Example: Allen-Cahn Equation}\label{example-allen-cahn-equation-1}

% \begin{figure}
%     \centering
%     \includegraphics{figures/AC_book.png}
%     \caption{Allen-Cahn equation: \textbf{Top:} Solution along
% with the location of the initial training snapshot at t=0.1 and the
% final prediction snapshot at t=0.9. \textbf{Bottom:} Initial training
% data and final prediction at the snapshots depicted by the white
% vertical lines in the top panel.}
% \end{figure}


% \subsection{Data-driven Discovery of Nonlinear Partial Differential
% Equations}\label{data-driven-discovery-of-nonlinear-partial-differential-equations}

% Parametrized and nonlinear partial differential equations of the general
% form

% \[u_t + \mathcal{N}[u;\lambda] = 0,\ x \in \Omega, \ t\in[0,T],\]

% where $u(t,x)$ denotes the latent (hidden) solution,
% $\mathcal{N}[\cdot;\lambda]$ is a nonlinear operator parametrized by
% $\lambda$, and $\Omega$ is a subset of $\mathbb{R}^D$.

% \paragraph{Continuous Time Models}\label{continuous-time-models}

% We define \[f(t,x)\] to be given by

% \[
% f := u_t + \mathcal{N}[u;\lambda],\label{eq:PDE_RHS}
% \]

% and proceed by approximating \[u(t,x)\] by a deep neural network. This
% assumption results in a physics informed neural network $f(t,x)$.

% \paragraph{Example: Navier-Stokes
% Equation}\label{example-navier-stokes-equation}

% Incompressible fluid flow as described by the ubiquitous Navier-Stokes
% equations. The Navier-Stokes equations in two dimensions (2D) given
% explicitly by

% \[
% \begin{array}{c}
% u_t + \lambda_1 (u u_x + v u_y) = -p_x + \lambda_2(u_{xx} + u_{yy}),\\
% v_t + \lambda_1 (u v_x + v v_y) = -p_y + \lambda_2(v_{xx} + v_{yy}),
% \end{array}
% \]

% where $u(t, x, y)$ denotes the $x$-component of the velocity field,
% $v(t, x, y)$ the $y$-component, and $p(t, x, y)$ the pressure. Here,
% $\lambda = (\lambda_1, \lambda_2)$ are the unknown parameters. Solutions
% to the Navier-Stokes equations are searched in the set of
% divergence-free functions; i.e.,

% \[
% u_x + v_y = 0.
% \]

% This extra equation is the continuity equation for incompressible fluids
% that describes the conservation of mass of the fluid. We make the
% assumption that

% \[
% u = \psi_y,\ \ \ v = -\psi_x,
% \]

% for some latent function $\psi(t,x,y)$. Under this assumption, the
% continuity equation will be automatically satisfied. Given noisy
% measurements

% \[
% \{t^i, x^i, y^i, u^i, v^i\}_{i=1}^{N}
% \]

% \paragraph{Example: Navier-Stokes Equation
% (contd.)}\label{example-navier-stokes-equation-contd.}

% We are interested in learning the parameters $\lambda$ as well as the
% pressure $p(t,x,y)$. We define $f(t,x,y)$ and $g(t,x,y)$ to be given by

% \[
% \begin{array}{c}
% f := u_t + \lambda_1 (u u_x + v u_y) + p_x - \lambda_2(u_{xx} + u_{yy}),\\
% g := v_t + \lambda_1 (u v_x + v v_y) + p_y - \lambda_2(v_{xx} + v_{yy}),
% \end{array}
% \]

% and proceed by jointly approximating \[\begin{bmatrix}
% \psi(t,x,y) & p(t,x,y)
% \end{bmatrix}\] using a single neural network with two outputs. This
% prior assumption results into a physics informed neural network
% \[\begin{bmatrix}
% f(t,x,y) & g(t,x,y)
% \end{bmatrix}\].

% The parameters $\lambda$ of the Navier-Stokes operator as well as the
% parameters of the neural networks
% $\begin{bmatrix} \psi(t,x,y) & p(t,x,y) \end{bmatrix}$ and
% $\begin{bmatrix} f(t,x,y) & g(t,x,y) \end{bmatrix}$ can be trained by
% minimizing the mean squared error loss

% \begin{equation}
% \begin{array}{rl}
% MSE :=& \frac{1}{N}\sum_{i=1}^{N} \left(|u(t^i,x^i,y^i) - u^i|^2 + |v(t^i,x^i,y^i) - v^i|^2\right) \\
%     +& \frac{1}{N}\sum_{i=1}^{N} \left(|f(t^i,x^i,y^i)|^2 + |g(t^i,x^i,y^i)|^2\right).
% \end{array}    
% \end{equation}



% Next: High predictive accuracy even when the temporal gap between the
% data snapshots is very large.

% \paragraph{Discrete Time Models}\label{discrete-time-models-1}

% We begin by employing the general form of
% \href{https://en.wikipedia.org/wiki/Runge--Kutta_methods}{Runge-Kutta}
% methods with \[q\] stages and obtain

% \[
% \begin{array}{ll}
% u^{n+c_i} = u^n - \Delta t \sum_{j=1}^q a_{ij} \mathcal{N}[u^{n+c_j};\lambda], \ \ i=1,\ldots,q,\\
% u^{n+1} = u^{n} - \Delta t \sum_{j=1}^q b_j \mathcal{N}[u^{n+c_j};\lambda].
% \end{array}
% \]

% Here, $u^{n+c_j}(x) = u(t^n + c_j \Delta t, x)$ for $j=1, \ldots, q$.
% This general form encapsulates both implicit and explicit time-stepping
% schemes, depending on the choice of the parameters $\{a_{ij},b_j,c_j\}$.
% The above equations can be equivalently expressed as

% \[
% \begin{array}{ll}
% u^{n} = u^n_i, \ \ i=1,\ldots,q,\\
% u^{n+1} = u^{n+1}_{i}, \ \ i=1,\ldots,q.
% \end{array}
% \]

% where

% \[
% \begin{array}{ll}
% u^n_i := u^{n+c_i} + \Delta t \sum_{j=1}^q a_{ij} \mathcal{N}[u^{n+c_j};\lambda], \ \ i=1,\ldots,q,\\
% u^{n+1}_{i} := u^{n+c_i} + \Delta t \sum_{j=1}^q (a_{ij} - b_j) \mathcal{N}[u^{n+c_j};\lambda], \ \ i=1,\ldots,q.
% \end{array}
% \]

% \paragraph{Discrete Time Models
% (contd.)}\label{discrete-time-models-contd.-1}

% We proceed by placing a multi-output neural network prior on

% \[
% \begin{bmatrix}
% u^{n+c_1}(x), \ldots, u^{n+c_q}(x)
% \end{bmatrix}.
% \]

% This prior assumption result in two physics informed neural networks

% \[
% \begin{bmatrix}
% u^{n}_1(x), \ldots, u^{n}_q(x), u^{n}_{q+1}(x)
% \end{bmatrix},
% \]

% and

% \[
% \begin{bmatrix}
% u^{n+1}_1(x), \ldots, u^{n+1}_q(x), u^{n+1}_{q+1}(x)
% \end{bmatrix}.
% \]

% \paragraph{Discrete Time Models
% (contd.)}\label{discrete-time-models-contd.-2}

% Given noisy measurements at two distinct temporal snapshots
% $\{\mathbf{x}^{n}, \mathbf{u}^{n}\}$ and
% $\{\mathbf{x}^{n+1}, \mathbf{u}^{n+1}\}$ of the system at times $t^{n}$
% and $t^{n+1}$, respectively, the shared parameters of the neural
% networks along with the parameters $\lambda$ of the differential
% operator can be trained by minimizing the sum of squared errors

% \[
% SSE = SSE_n + SSE_{n+1},
% \]

% where

% \[
% SSE_n := \sum_{j=1}^q \sum_{i=1}^{N_n} |u^n_j(x^{n,i}) - u^{n,i}|^2,
% \]

% and

% \[
% SSE_{n+1} := \sum_{j=1}^q \sum_{i=1}^{N_{n+1}} |u^{n+1}_j(x^{n+1,i}) - u^{n+1,i}|^2.
% \]

% Here, $\mathbf{x}^n = \left\{x^{n,i}\right\}_{i=1}^{N_n}$,
% $\mathbf{u}^n = \left\{u^{n,i}\right\}_{i=1}^{N_n}$,
% $\mathbf{x}^{n+1} = \left\{x^{n+1,i}\right\}_{i=1}^{N_{n+1}}$, and
% $\mathbf{u}^{n+1} = \left\{u^{n+1,i}\right\}_{i=1}^{N_{n+1}}$.

% \paragraph{Example: Korteweg--de Vries
% Equation}\label{example-kortewegde-vries-equation}

% The final example aims to highlight the ability of the proposed
% framework to handle governing partial differential equations involving
% higher order derivatives. Here, we consider a mathematical model of
% waves on shallow water surfaces; the Korteweg-de Vries (KdV) equation.
% The KdV equation reads as

% \[
% u_t + \lambda_1 u u_x + \lambda_2 u_{xxx} = 0,
% \]

% with $(\lambda_1, \lambda_2)$ being the unknown parameters. For the KdV
% equation, the nonlinear operator is given by

% \[
% \mathcal{N}[u^{n+c_j}] = \lambda_1 u^{n+c_j} u^{n+c_j}_x - \lambda_2 u^{n+c_j}_{xxx}
% \]

% and the shared parameters of the neural networks along with the
% parameters \[\lambda = (\lambda_1, \lambda_2)\] of the KdV equation can
% be learned by minimizing the sum of squared errors given above.

% \paragraph{Example: Korteweg--de Vries Equation
% (contd.)}\label{example-kortewegde-vries-equation-contd.}





% \subsection{Summary}\label{summary}

% \begin{itemize}
% \item
%   PINNs provide a reqularization framework based on the (partially)
%   known physics.
% \item
%   Spartially Low-Data Regime: Continuous-Time Models
% \item
%   Temporally Low-Data Regime: Discrete-Time Models (RK Method)
% \end{itemize}

% \subsection{Open Questions}\label{open-questions}

% How deep/wide should the neural network be? How much data is really
% needed?

% Why does the algorithm converge to unique values for the parameters of
% the differential operators?

% Does the network suffer from vanishing gradients for deeper
% architectures and higher order differential operators?

% Could this be mitigated by using different activation functions?

% Can we improve on initializing the network weights or normalizing the
% data?

% Are the mean square error and the sum of squared errors the appropriate
% loss functions?

% Why are these methods seemingly so robust to noise in the data?

% How can we quantify the uncertainty associated with our predictions?


\end{document}
